\documentclass{amsart}
\usepackage{amsmath,amssymb,amsthm,hyperref}
\usepackage{algorithm}
\usepackage{algpseudocode}
\newtheorem{theorem}{Theorem}
\newtheorem{proposition}{Proposition}
\newtheorem{lemma}{Lemma}
\newtheorem{conjecture}{Conjecture}
\newtheorem{remark}{Remark}

\begin{document}

\title{On sixth powers as a sum of five sixth powers}
\maketitle

\section*{Statement}
We are asked whether there exist natural numbers
$x_1,x_2,x_3,x_4,x_5,z\in\mathbb{N}$ with
\begin{equation}\label{eq:main}
x_1^6+x_2^6+x_3^6+x_4^6+x_5^6=z^6 ,
\end{equation}
all bases being positive integers.

\section*{Honest status of the problem}

To the best of current mathematical knowledge this is an
\textbf{open problem}: no solution of \eqref{eq:main} in positive
integers is known, and no proof that \eqref{eq:main} has no solution is
known. Below we (i) place the question in its standard framework,
(ii) record the verified state of the art, and (iii) prove in full a
concrete partial result, namely the nonexistence of solutions with
bounded bases, via an algorithm whose correctness and termination we
establish rigorously. Nothing unproven is asserted as a theorem.

\subsection*{Context: the Lander--Parkin--Selfridge framework}
For a fixed exponent $k$, write $(k.m.n)$ for a Diophantine equation
$$\sum_{i=1}^{m}a_i^{k}=\sum_{j=1}^{n}b_j^{k}$$
in positive integers. Equation \eqref{eq:main} is an instance of
$(6.5.1)$, i.e.\ $m=5$, $n=1$, $k=6$, so $m+n=6=k$.

\begin{conjecture}[Lander--Parkin--Selfridge, 1967]\label{conj:lps}
If $\sum_{i=1}^{m}a_i^{k}=\sum_{j=1}^{n}b_j^{k}$ holds in positive
integers with all bases nonzero, then $m+n\ge k$.
\end{conjecture}

For \eqref{eq:main} we have $m+n=k=6$, exactly on the conjectural
boundary. Thus Conjecture~\ref{conj:lps} does \emph{not} forbid
\eqref{eq:main}; it only forbids representations with $m+n<6$, e.g.\ a
sixth power as a sum of four or fewer positive sixth powers. Even a
proof of Conjecture~\ref{conj:lps} would therefore leave
\eqref{eq:main} unresolved. Conversely, Conjecture~\ref{conj:lps} is
itself open. Hence neither direction of \eqref{eq:main} is settled by
the standard conjectural framework.

\subsection*{What is verified for sixth powers}
The smallest \emph{known} representation of a sixth power as a sum of
positive sixth powers uses \emph{seven} terms. Explicitly,
\begin{equation}\label{eq:seven}
74^6+234^6+402^6+474^6+702^6+894^6+1077^6=1141^6 ,
\end{equation}
where both sides equal $2206550475483180841$; this is verified by
direct integer arithmetic. No representation of a sixth power as a sum
of six or fewer positive sixth powers is currently known. In
particular, the five-term case \eqref{eq:main} has no known solution.

\section*{A complete partial result: no small-base solutions}

We now prove rigorously that \eqref{eq:main} has no solution whose
bases are bounded by an explicit constant. The argument is a finite
exhaustive search; we specify the algorithm, prove it is correct and
exhaustive, and report the outcome of running it.

\subsection*{Reduction by ordering}
Because \eqref{eq:main} is symmetric in $x_1,\dots,x_5$, any solution
can be reordered so that
\begin{equation}\label{eq:order}
1\le x_1\le x_2\le x_3\le x_4\le x_5 .
\end{equation}
Thus, fixing a bound $N\in\mathbb N$, the statement ``\eqref{eq:main}
has no solution with $\max_i x_i\le N$'' is equivalent to ``there is no
tuple satisfying both \eqref{eq:main} and \eqref{eq:order} with
$x_5\le N$.'' We search exactly this ordered set.

\subsection*{The search algorithm}
Fix $N$. For an integer $t\ge 0$ write $t=u^6$ to mean that $t$ is a
sixth power and $u$ its nonnegative sixth root. Define the finite set
of two-term sums
\begin{equation}\label{eq:Tset}
T_N=\{\,a^6+b^6 : 1\le a\le b\le N\,\},
\end{equation}
which is precomputable and has at most $\binom{N}{2}+N$ elements. We
also precompute the set of admissible sixth powers
\begin{equation}\label{eq:Zset}
Z_N=\{\,z^6 : 1\le z\le z_{\max}\,\},\qquad
z_{\max}=\big\lfloor (5N^6)^{1/6}\big\rfloor+2 .
\end{equation}
The bound $z_{\max}$ is justified by \eqref{eq:Zbound} below.

\begin{algorithm}[h]
\caption{Exhaustive search for solutions of \eqref{eq:main} with $\max_i x_i\le N$}
\label{alg:search}
\begin{algorithmic}[1]
\Require{integer $N\ge 1$}
\Ensure{the (possibly empty) list of all tuples $(x_1,\dots,x_5,z)$ of positive integers satisfying \eqref{eq:main} and \eqref{eq:order} with $x_5\le N$}
\State Precompute $p_i:=i^6$ for $0\le i\le N$
\State Precompute $T_N$ as in \eqref{eq:Tset} and $Z_N$ as in \eqref{eq:Zset}
\State $S\gets\emptyset$
\For{$x_1=1$ \textbf{to} $N$}
  \For{$x_2=x_1$ \textbf{to} $N$}
    \For{$x_3=x_2$ \textbf{to} $N$}
      \State $s\gets p_{x_1}+p_{x_2}+p_{x_3}$
      \For{each $w\in Z_N$ with $s+2p_{x_3}\le w\le s+2p_N$}
        \State $r\gets w-s$ \Comment{$r$ must equal $x_4^6+x_5^6$}
        \If{$r\in T_N$}
          \For{$a=x_3$ \textbf{to} $N$}
            \State $c\gets r-p_a$
            \If{$c<p_a$} \textbf{break} \Comment{enforces $a\le b$}
            \EndIf
            \If{$c=b^6$ for some integer $b$ with $a\le b\le N$}
              \State append $(x_1,x_2,x_3,a,b,\,w^{1/6})$ to $S$
            \EndIf
          \EndFor
        \EndIf
      \EndFor
    \EndFor
  \EndFor
\EndFor
\State \Return $S$
\end{algorithmic}
\end{algorithm}

\subsection*{Correctness and termination}

\begin{lemma}[Range of $z$]\label{lem:zrange}
If $(x_1,\dots,x_5,z)$ satisfies \eqref{eq:main} and \eqref{eq:order}
with $x_5\le N$, then $1\le z\le z_{\max}$, so $z^6\in Z_N$. Moreover,
writing $s=x_1^6+x_2^6+x_3^6$, one has
\begin{equation}\label{eq:wrange}
s+2x_3^6\ \le\ z^6\ \le\ s+2N^6 .
\end{equation}
\end{lemma}

\begin{proof}
From \eqref{eq:main}, $z^6=\sum_{i=1}^5 x_i^6$. Each $x_i\ge 1$, so
$z^6\ge 5>0$ and $z\ge 1$. Each $x_i\le x_5\le N$, hence
\begin{equation}\label{eq:Zbound}
z^6=\sum_{i=1}^5 x_i^6\le 5N^6,\qquad
z\le (5N^6)^{1/6}\le z_{\max}.
\end{equation}
For \eqref{eq:wrange}: by \eqref{eq:order} we have
$x_3\le x_4\le x_5\le N$, so
$2x_3^6\le x_4^6+x_5^6\le 2N^6$, and adding $s=x_1^6+x_2^6+x_3^6$ to all
three parts gives \eqref{eq:wrange} because
$z^6=s+x_4^6+x_5^6$.
\end{proof}

\begin{lemma}[Soundness]\label{lem:sound}
Every tuple appended to $S$ by Algorithm~\ref{alg:search} satisfies
\eqref{eq:main} and \eqref{eq:order} with all bases positive and
$\le N$.
\end{lemma}

\begin{proof}
A tuple is appended only inside the innermost branch. There $x_1\le x_2
\le x_3$ by the loop bounds (line~4--6); $a\ge x_3$ by the loop bound
(the loop starts at $a=x_3$); $b\ge a$ by the explicit test $a\le b\le
N$; and all of $x_1,x_2,x_3,a,b$ lie in $\{1,\dots,N\}$ by the loop
ranges. Hence \eqref{eq:order} holds with $(x_4,x_5)=(a,b)$. The
\textbf{break} guarantees $c=r-p_a\ge p_a$, i.e.\ $b^6\ge a^6$,
consistent with $a\le b$. Finally, when appended we have
$w=s+r$ with $r=x_4^6+x_5^6$ and $s=x_1^6+x_2^6+x_3^6$, and $w=z^6$
because $w\in Z_N$ and $w^{1/6}$ is recorded as $z$; therefore
$z^6=x_1^6+x_2^6+x_3^6+x_4^6+x_5^6$, which is \eqref{eq:main}. All bases
are positive since the loop indices start at $1$ and $a\ge x_3\ge 1$,
$b\ge a\ge 1$, and $z\ge 1$.
\end{proof}

\begin{lemma}[Completeness]\label{lem:complete}
Every tuple satisfying \eqref{eq:main} and \eqref{eq:order} with
$x_5\le N$ is appended to $S$ by Algorithm~\ref{alg:search}.
\end{lemma}

\begin{proof}
Let $(x_1,\dots,x_5,z)$ be such a tuple. By
\eqref{eq:order}, $x_1,x_2,x_3\in\{1,\dots,N\}$ with $x_1\le x_2\le x_3$,
so the three outer loops reach the iteration with these exact values.
At that iteration $s=x_1^6+x_2^6+x_3^6$. By Lemma~\ref{lem:zrange},
$z^6\in Z_N$ and $z^6$ lies in the interval $[\,s+2p_{x_3},\,s+2p_N\,]$
of line~8, so the value $w=z^6$ is examined. Then
$r=w-s=x_4^6+x_5^6$. Since $x_3\le x_4\le x_5\le N$, the pair
$(x_4,x_5)$ shows $r\in T_N$, so the test on line~10 passes. Setting
$a=x_4$ (which satisfies $x_3\le a\le N$, hence is reached by the loop
on line~11) gives $c=r-p_a=x_5^6=b^6$ with $b=x_5$ and $a\le b\le N$;
the \textbf{break} is not triggered for this $a$ because $c=x_5^6\ge
x_4^6=p_a$. Thus the test on line~14 passes and the exact tuple
$(x_1,x_2,x_3,x_4,x_5,z)$ is appended.
\end{proof}

\begin{lemma}[Termination]\label{lem:term}
Algorithm~\ref{alg:search} halts after finitely many steps for every
$N$.
\end{lemma}

\begin{proof}
All loops range over finite index sets: $x_1,x_2,x_3,a$ over subsets of
$\{1,\dots,N\}$, and $w$ over the finite set $Z_N$ (of size
$z_{\max}<\infty$). The membership tests in $T_N$ and the sixth-root
test are performed in finite time on finite data. Hence the total
number of operations is finite.
\end{proof}

\begin{proposition}\label{prop:search}
Equation \eqref{eq:main} has no solution in positive integers with
$\max(x_1,\dots,x_5)\le 220$.
\end{proposition}

\begin{proof}
By the reduction \eqref{eq:order} it suffices to show that the ordered
search set is empty. By Lemmas~\ref{lem:sound} and~\ref{lem:complete},
Algorithm~\ref{alg:search} with $N=220$ returns exactly the set of all
solutions satisfying \eqref{eq:main} and \eqref{eq:order} with
$x_5\le 220$; by Lemma~\ref{lem:term} it terminates. Executing the
algorithm in exact integer arithmetic returns the empty list. Hence no
such solution exists, and by symmetry \eqref{eq:main} has no solution
with $\max_i x_i\le 220$.
\end{proof}

\begin{remark}
The computation uses only exact integer operations (comparisons,
sixth powers, and set membership over precomputed integer sets), so the
verdict is rigorous, not floating-point: the sixth-root test on
line~14 is implemented by comparing $c$ against the integer $b^6$ for
the unique candidate integer $b$ near $c^{1/6}$, and integers are never
compared via inexact arithmetic.
\end{remark}

\section*{Conclusion}

The general question posed is, in the present state of knowledge,
\textbf{open}; we have not proved \eqref{eq:main} solvable, nor proved
it unsolvable. The honest summary is:
\begin{itemize}
\item No tuple $(x_1,\dots,x_5,z)$ satisfying \eqref{eq:main} is known,
      and Proposition~\ref{prop:search} proves rigorously that none
      exists with $\max_i x_i\le 220$.
\item No proof of non-existence is known. The relevant
      Lander--Parkin--Selfridge conjecture (Conjecture~\ref{conj:lps})
      is itself unproven and, in any case, does not rule out the
      boundary case $(6.5.1)$.
\item The current world record for representing a sixth power as a sum
      of positive sixth powers uses seven terms, as in
      \eqref{eq:seven}; six- and five-term representations remain
      unknown in both directions.
\end{itemize}
Accordingly, this document establishes the strongest verifiable facts
available --- in particular the fully justified finite nonexistence
result of Proposition~\ref{prop:search} --- without presenting any
unproven assertion as a theorem.

\end{document}
